\chapter{COSHH Assessment for Construction Chemicals}
\label{chap:COSHH Assessment for Construction Chemicals}

%---------------------------- SECTION ----------------------------
\section{Why this assessment matters in construction}

Construction operatives are exposed to a wider range of hazardous substances than workers in almost any other industry. Cement and concrete, epoxy resins, isocyanates, solvents, bitumen, wood preservatives, lead-containing paints, hydraulic oils, and a multitude of proprietary adhesives, coatings, and sealants are used across the trades on a daily basis. Many of these substances cause acute injury on first contact; many cause chronic disease that develops over years of repeated low-level exposure and manifests only when it is too late to prevent or reverse. The HSE estimates that approximately 12,000 deaths per year in the UK are attributable to occupational cancers, with construction trades represented disproportionately in this total. The COSHH Regulations 2002 provide the framework to identify, assess, and control chemical hazards before they cause harm.

Cement dermatitis provides a telling example of the gap between awareness and control in the construction industry. Chrome VI (chromate) in cement is one of the most well-understood occupational skin sensitisers in any sector. Regulatory action to reduce chrome VI levels in cement was implemented across the EU in 2005. Construction operatives have been told for decades that wet cement causes dermatitis. And yet, wet work dermatitis --- particularly chrome VI cement dermatitis and irritant contact dermatitis from high-alkalinity cement --- remains one of the most prevalent occupational skin diseases in UK construction. The explanation is not ignorance of the hazard; it is the failure of COSHH assessments to translate the identification of the hazard into effective, sustained, and supervised control.

The Control of Substances Hazardous to Health Regulations 2002 define the eight-step COSHH assessment process: identify substances used; decide if they are hazardous; determine how, where, and how often they are used; estimate likelihood and level of exposure; decide whether controls are adequate; record findings; review and update; carry out surveillance where appropriate. This process is well understood by safety professionals but is routinely applied in construction as a pro-forma exercise in which the Safety Data Sheet (SDS) for each product is reviewed and the conclusion is reached that exposure can be managed with gloves and PPE, without the detailed consideration of engineering controls, substitution, and work method redesign that the Regulations actually require.

The introduction of the reach regulation and updated CLP classification system has added complexity to the management of construction chemicals, with frequent reformulations of products that were assessed previously now requiring reassessment. The assessor must ensure that Safety Data Sheets are current, that products have not been reformulated since the last assessment, and that Workplace Exposure Limits (WELs) from EH40 are applied to all substances with assigned WELs.

\barquote{``A COSHH assessment that concludes with `wear gloves and PPE' without considering engineering controls has not followed the hierarchy; it has documented an exposure that could have been reduced.''}

\example{Site Reality}{A tiling subcontractor is working on a bathroom pod installation contract in a hotel development. The grout used is a two-part epoxy grout containing an epoxy resin component and a hardener. The COSHH assessment prepared by the contractor specifies nitrile gloves as the control. Three months into the contract, two tilers develop skin sensitisation to epoxy resin. Medical investigation confirms occupational contact dermatitis with complete sensitisation --- both operatives will never be able to work with epoxy products again. The investigation finds that the nitrile gloves specified in the COSHH assessment were standard examination gloves with a thickness insufficient to prevent epoxy penetration over the extended contact periods involved in grout application; the assessment had not addressed the contact duration or the need for thicker chemical-resistant gloves; and no biological monitoring or health surveillance had been implemented despite the known sensitisation risk of epoxy resin products.}

%---------------------------- SECTION ----------------------------
\section{Legal and professional framework}

The Control of Substances Hazardous to Health Regulations 2002 (COSHH) are the primary legislative instrument for chemical hazard management in construction. They apply to any substance that presents a health hazard through inhalation, skin absorption, ingestion, or injection, including dusts (covered in Chapter~\ref{chap:Silica, Wood Dust and Respiratory Hazard Risk Assessment}), biological agents, and chemicals. The COSHH Regulations impose a hierarchy of control identical in structure to the general hierarchy: eliminate, substitute, engineer, administrate, PPE. EH40 (Workplace Exposure Limits) provides the current list of WELs and occupational exposure limits for substances in the UK workplace, and is updated periodically to reflect new toxicological data.

The Classification, Labelling and Packaging (CLP) Regulation and the Registration, Evaluation, Authorisation and Restriction of Chemicals (REACH) Regulation govern how chemical substances are classified, labelled, and their use communicated via Safety Data Sheets. REACH also restricts the use of substances of very high concern, and construction contractors must be aware of SVHC restrictions that may affect products they procure and use. CDM 2015 requires designers to specify less hazardous materials where reasonably practicable and to provide pre-construction information about hazardous materials present in the existing structure.

\paragraph{Key frameworks relevant to this chapter include: COSHH Regulations 2002; EH40 Workplace Exposure Limits; CLP Regulation; REACH Regulation; CDM 2015.}

\begin{itemize}
  \bitem{COSHH Regulations 2002}{Impose the eight-step COSHH assessment process; require adequate control of exposure to hazardous substances; require maintenance, examination, and testing of control measures; require health surveillance where there is a reasonable likelihood of an identifiable disease associated with the substance; require that employees are informed, instructed, and trained regarding the hazards of substances they use.}
  \bitem{EH40 Workplace Exposure Limits}{Provides the WELs for all substances with assigned limits in UK workplaces; WELs must be observed as legal maximum limits in the case of substances with a maximum exposure limit designation; compliance with WELs must be verified through exposure monitoring where appropriate.}
  \bitem{CLP Regulation and Safety Data Sheets}{Requires that all hazardous chemicals are properly classified, labelled, and accompanied by a Safety Data Sheet in accordance with the standard 16-section format; SDSs must be current (version checked against product formulation), available to workers, and used as the foundation of the COSHH assessment.}
  \bitem{CDM 2015}{Requires designers to eliminate the use of hazardous substances through design choices where reasonably practicable (for example, specifying water-based rather than solvent-based products); requires that pre-construction information identifies hazardous materials in the existing structure (asbestos, lead paint, contaminated ground) before work commences.}
\end{itemize}

%---------------------------- SECTION ----------------------------
\section{Conducting the assessment using the five-step method}

\subsection{Step 1 --- Identify hazards}
\paragraph{COSHH hazards in construction include: cement and concrete dust causing cement dermatitis (chrome VI sensitisation) and respiratory irritation; wet cement causing severe alkali burns to skin and eyes; epoxy resins and hardeners causing skin sensitisation and occupational contact dermatitis; isocyanates in spray foam insulation causing occupational asthma (a prescribed occupational disease under RIDDOR); solvent-based adhesives, coatings, and sealants causing narcosis, respiratory irritation, and chronic neurological damage from long-term exposure; bitumen fumes from hot bitumen application causing respiratory irritation and skin exposure to polycyclic aromatic hydrocarbons (possible carcinogen); lead paint on existing surfaces encountered during sanding, grinding, or torch-cutting; hydraulic oils and cutting fluids causing dermatitis and occupational dermatological conditions; and wood preservatives containing biocides and copper compounds causing skin and respiratory sensitisation.}

\subsection{Step 2 --- Decide who or what is affected}
\paragraph{Persons at risk from construction chemical hazards include: all operatives directly handling or applying the substance; operatives working in the immediate vicinity where airborne contaminants are generated (spray application, grinding, sanding); workers in confined spaces where solvent vapours can accumulate; apprentices and young workers who may have no prior awareness of chemical sensitisation risk and who, if sensitised, will carry that sensitisation for the rest of their working lives; and members of the public adjacent to applications generating airborne contaminant plumes (spray foam operations, solvent-based coating applications). Particular attention must be given to skin conditions: even short-term exposure to sensitisers such as epoxy resin and chrome VI cement can produce lifelong sensitisation that ends the affected person's ability to work in trades using those substances.}

\subsection{Step 3 --- Evaluate likelihood and consequence}
\paragraph{For substances with chronic health effects (carcinogens, sensitisers, reproductive hazards), the consequence of uncontrolled exposure is career-ending or life-limiting disease; the inherent risk must therefore be rated High to Very High (15--25 on a 5×5 matrix) where exposure is routine and controls are absent. For acute hazards (alkali burns from cement, eye injury from solvent splash), the consequence is potentially severe and the likelihood is High on active work sites. The COSHH hierarchy, when properly applied, can reduce exposure to below the WEL and reduce residual risk to Low; the reduction must be demonstrated through exposure monitoring where WEL-applicable substances are involved.}

\subsection{Step 4 --- Select and implement controls}
\paragraph{COSHH controls must follow the hierarchy: first, eliminate (can a product be removed from the specification entirely? Can an activity that generates exposure be redesigned?); second, substitute (can a less hazardous product be used? Water-based rather than solvent-based; low-chrome cement already reduced by Regulation; non-isocyanate foam rather than MDI-based spray foam?); third, engineer (local exhaust ventilation (LEV) for dust generation, solvent mixing, and spray application; on-tool extraction; enclosed application systems; dilution ventilation); fourth, administrative (reduce the number of operatives exposed; minimise exposure duration; implement LEV maintenance and testing schedule; provide health surveillance; train operatives in correct technique and product hazards); fifth, PPE as the final layer (chemical-resistant gloves of adequate thickness and type; eye protection; respiratory protective equipment appropriate to the contaminant).}

\subsection{Step 5 --- Record and review}
\paragraph{The COSHH assessment must be reviewed whenever a product is reformulated, substituted, or replaced; whenever a new hazardous substance is introduced onto the site; whenever a health surveillance finding indicates that exposure is causing harm; and at least annually for ongoing activities. Product Safety Data Sheets must be checked for currency at the time of each review. Exposure monitoring results, health surveillance records, and LEV test records must all be retained as part of the COSHH file.}

\example{Five-Step Application}{A heating and plumbing contractor is installing underfloor heating in a new residential development using a solvent-based adhesive to bond the insulation boards to the concrete subfloor. Step 1 identifies: solvent vapour (MEK and toluene-based) accumulating in the ground-floor slab area, which is enclosed during winter working. Step 2 identifies: two plumbers applying adhesive; electrician working in the adjacent area. Step 3 rates inherent solvent vapour exposure risk as High (16/25) in the enclosed environment. Step 4 implements: water-based equivalent adhesive specified as substitution after supplier consultation; enclosed area ventilated with temporary forced ventilation before and during application; WEL monitoring for residual toluene confirms levels below WEL; nitrile chemical-resistant gloves and RPE half-mask with organic vapour cartridges specified for application only. Residual risk rated Low (6/25). Step 5: SDS review before each delivery batch to confirm formulation unchanged.}

%---------------------------- SECTION ----------------------------
\section{Applying the hierarchy of controls}

\paragraph{The COSHH hierarchy is directly aligned with the general control hierarchy and must be applied in strict sequence; PPE is the least reliable control and must not be specified as the primary means of achieving adequate control unless the earlier controls in the hierarchy have been genuinely considered and rejected.}

\begin{itemize}
  \bitem{Elimination}{Remove the hazardous substance from the specification: specify water-based coatings; choose mechanical fixing rather than chemical anchors where structural performance permits; specify pre-finished or factory-applied products that eliminate the need for on-site coating application.}
  \bitem{Substitution}{Replace the higher-hazard product with a lower-hazard alternative: substitute solvent-based adhesives with water-based equivalents; substitute two-component isocyanate foam with water-expanding polyurethane alternatives; specify low-chrome cement to reduce chrome VI exposure.}
  \bitem{Engineering controls}{Local exhaust ventilation at the point of generation; on-tool extraction for grinding and cutting; enclosed application systems for spray products; dilution ventilation (minimum 10 air changes per hour) in enclosed spaces where solvents are used; controlled dispensing equipment to minimise spill risk.}
  \bitem{Administrative controls}{Restrict access to areas of high chemical exposure to those directly involved in the activity; minimise the number of operatives exposed and the duration of exposure; implement health surveillance for sensitiser-associated substances (epoxy, isocyanates, chrome VI); maintain LEV inspection and maintenance records; train all operatives in COSHH awareness, product hazards, and correct use of controls; ensure current SDS is available at the point of use.}
  \bitem{PPE}{Chemical-resistant gloves appropriate to the substance (nitrile for most solvents; laminated barrier gloves for epoxy resins; check SDS for specific glove material recommendations); chemical splash goggles; RPE appropriate to the substance and exposure level (disposable FFP3 for inorganic dust; half-mask with appropriate cartridge for organic vapours; powered air purifying respirator (PAPR) for spray operations).}
\end{itemize}

%---------------------------- SECTION ----------------------------
\section{Cadence of assessment and review triggers}

\paragraph{COSHH assessments must be reviewed on a regular cycle and in response to specific trigger events; the assessment cycle is driven by both the rate of change in chemical products and the known latency of occupational disease.}

\begin{itemize}
  \bitem{When a new product is introduced}{Any new chemical product brought onto site must have a current SDS reviewed and a COSHH assessment completed or updated before it is used; the introduction of a new product must not be permitted until this is complete.}
  \bitem{When a product formulation changes}{Product reformulations may change the hazard classification, WEL applicability, or recommended controls without a change in the product name; the SDS must be checked with each new delivery batch to confirm no reformulation has occurred.}
  \bitem{When a health surveillance finding is adverse}{Any finding from health surveillance (skin assessment, lung function test, biological monitoring) that indicates adverse effect must trigger an immediate COSHH assessment review and a review of the effectiveness of current controls.}
  \bitem{When working conditions change}{A change from outdoor to indoor working, from dry to wet weather, or from ventilated to enclosed workspace changes the exposure conditions and may invalidate the assumptions of the original assessment.}
  \bitem{Minimum annually}{All COSHH assessments for ongoing activities must be reviewed at least annually to confirm that products, working methods, and controls remain as described in the assessment.}
\end{itemize}

%---------------------------- SECTION ----------------------------
\section{Common failure modes}

\begin{itemize}
  \bitem{SDS relied upon without assessment}{The SDS is treated as the COSHH assessment. Safety data sheets are hazard communication documents, not risk assessments; they describe the hazard, not the risk in the specific workplace context.}
  \bitem{PPE specified as primary control without hierarchy consideration}{Every COSHH assessment on the project specifies gloves and RPE as the control measures, without demonstrating that engineering controls were considered and found impracticable. This is a regulatory failure and a missed opportunity to protect workers more effectively.}
  \bitem{Health surveillance not implemented for sensitising substances}{Isocyanates, epoxy resins, and chrome VI cement are all recognised occupational sensitisers for which health surveillance is indicated under COSHH. Health surveillance is routinely absent from construction COSHH frameworks.}
  \bitem{SDSs not current}{Assessments are based on SDSs that are several years old and do not reflect product reformulations. The assessor does not know that the product has been reformulated because the product name has not changed.}
  \bitem{LEV not tested or maintained}{Local exhaust ventilation systems installed as a control measure are not maintained or tested. An untested LEV system may provide a false assurance of control while actually providing negligible exposure reduction.}
  \bitem{Inadequate glove selection}{Generic ``chemical-resistant gloves'' are specified without reference to the specific product, contact duration, or breakthrough time data in the SDS. Gloves that are adequate for incidental contact may fail within minutes during prolonged direct handling.}
\end{itemize}

\example{Failure Pattern}{A flooring contractor is applying a two-component polyurethane floor coating to a warehouse. The COSHH assessment specifies nitrile gloves, safety glasses, and ``adequate ventilation''. The warehouse has no natural ventilation and temporary forced ventilation is not installed. The solvent component of the coating generates vapour levels well above the WEL; the isocyanate hardener component generates respiratory sensitiser levels that, over the two-week application period, sensitise one of the two applicators to isocyanates. The sensitised operative is diagnosed with occupational asthma twelve months later; he will never again be able to work in any environment where isocyanates are present. The COSHH assessment had identified isocyanates as a hazard; it had not considered health surveillance, adequate ventilation as an engineering control, or the use of supplied-air breathing apparatus for spray application of isocyanate products.}

%---------------------------- CHAPTER SUMMARY ----------------------------
\section{Chapter Summary}

\begin{table}[H]
\centering
\begin{tabularx}{\textwidth}{>{\raggedright\arraybackslash}p{3.2cm}X}
\toprule
\textbf{Assessment Element} & \textbf{Operational Requirement} \\
\midrule
Legal/Professional Basis & COSHH Regulations 2002; EH40; CLP Regulation; CDM 2015. Eight-step COSHH process with hierarchy from elimination to PPE must be demonstrated. \\
Method & Apply the full five-step process and document inherent/residual risk. \\
Controls & Substitution of less hazardous products; LEV and engineering controls before PPE; health surveillance for sensitising substances; current SDS at point of use. \\
Cadence & When new products are introduced; when formulations change; when health surveillance findings are adverse; minimum annually. \\
Assurance & Use audit, incident learning, and governance reporting to verify effectiveness. \\
\bottomrule
\end{tabularx}
\end{table}

\begin{itemize}
  \bitem{Key takeaway 1}{A COSHH assessment that begins and ends with PPE selection has not followed the hierarchy required by the Regulations; substitution, engineering controls, and administrative measures must be genuinely considered before PPE is specified.}
  \bitem{Key takeaway 2}{Skin sensitisation from epoxy resin, chrome VI cement, and isocyanates is a career-ending consequence that occurs after exposures that would not be classified as serious incidents at the time; health surveillance is the detection mechanism that must be built into the COSHH management system.}
  \bitem{Key takeaway 3}{Safety Data Sheets must be current and reviewed with each new delivery batch; product reformulations change hazard classifications and recommended controls without a change in product name.}
  \bitem{Key takeaway 4}{LEV systems are only effective if maintained and tested; an untested LEV installation is not a demonstrated control and must not be recorded as such in the COSHH assessment.}
\end{itemize}

\barquote{In Chapter~\ref{chap:Asbestos Management Risk Assessment}, we turn from this assessment to the next domain in the Prudentia framework.}
